%%%%%%%%%%%%%%%%%%%%%%%%%%%%%%%%%%%%%%%%%
% Jacobs Landscape Poster
% LaTeX Template
% Version 1.0 (29/03/13)
%
% Created by:
% Computational Physics and Biophysics Group, Jacobs University
% https://teamwork.jacobs-university.de:8443/confluence/display/CoPandBiG/LaTeX+Poster
%
% Further modified by:
% Nathaniel Johnston (nathaniel@njohnston.ca)
%
% This template has been downloaded from:
% http://www.LaTeXTemplates.com
%
% License:
% CC BY-NC-SA 3.0 (http://creativecommons.org/licenses/by-nc-sa/3.0/)
%
%%%%%%%%%%%%%%%%%%%%%%%%%%%%%%%%%%%%%%%%%

%----------------------------------------------------------------------------------------
%	PACKAGES AND OTHER DOCUMENT CONFIGURATIONS
%----------------------------------------------------------------------------------------

\documentclass[final]{beamer}

\usepackage[scale=1.24]{beamerposter} % Use the beamerposter package for laying out the poster
\usepackage{pgfplots}
\usepackage{float}
\usepackage{onimage}
\usepackage{etex}
\usetikzlibrary{arrows}

\usetheme{confposter} % Use the confposter theme supplied with this template

\setbeamercolor{block title}{fg=ngreen,bg=white} % Colors of the block titles
\setbeamercolor{block body}{fg=black,bg=white} % Colors of the body of blocks
\setbeamercolor{block alerted title}{fg=white,bg=dblue!70} % Colors of the highlighted block titles
\setbeamercolor{block alerted body}{fg=black,bg=dblue!10} % Colors of the body of highlighted blocks
% Many more colors are available for use in beamerthemeconfposter.sty

%-----------------------------------------------------------
% Define the column widths and overall poster size
% To set effective sepwid, onecolwid and twocolwid values, first choose how many columns you want and how much separation you want between columns
% In this template, the separation width chosen is 0.024 of the paper width and a 4-column layout
% onecolwid should therefore be (1-(# of columns+1)*sepwid)/# of columns e.g. (1-(4+1)*0.024)/4 = 0.22
% Set twocolwid to be (2*onecolwid)+sepwid = 0.464
% Set threecolwid to be (3*onecolwid)+2*sepwid = 0.708

\newlength{\sepwid}
\newlength{\onecolwid}
\newlength{\twocolwid}
\newlength{\threecolwid}
\setlength{\paperwidth}{48in} % A0 width: 46.8in
\setlength{\paperheight}{36in} % A0 height: 33.1in
\setlength{\sepwid}{0.024\paperwidth} % Separation width (white space) between columns
\setlength{\onecolwid}{0.22\paperwidth} % Width of one column
\setlength{\twocolwid}{0.464\paperwidth} % Width of two columns
\setlength{\threecolwid}{0.708\paperwidth} % Width of three columns
\setlength{\topmargin}{-0.5in} % Reduce the top margin size
%-----------------------------------------------------------

\usepackage{graphicx}  % Required for including images

\usepackage{booktabs} % Top and bottom rules for tables

% Macros for arrows between images
% based on http://tex.stackexchange.com/questions/7925/big-arrows-between-images
%----------------------------------------------------------------------------------------
\newcommand*{\vimage}[1]{\vcenter{\hbox{#1}}}
\newcommand*{\vpointer}[1][\ \triangleright\ ]{\vcenter{\hbox{$#1$}}}
\usepackage{mathtools} % for xRightarrow

\tikzset{MyCaption/.style={white, fill=black, opacity=0.95, text opacity=1.0,
rounded corners=0.5ex, anchor=center}}

% Tikz stuff for the cool diagrams
%----------------------------------------------------------------------------------------
\usetikzlibrary{spy}

%----------------------------------------------------------------------------------------
%	TITLE SECTION
%----------------------------------------------------------------------------------------

\title{Finding Eigenvalues using Automatic Differentiation} % Poster title

\author{Alex Hirzel, Peter Solfest, and Ryan Bruner} % Author(s)

\institute{Michigan Technological University} % Institution(s)

%----------------------------------------------------------------------------------------

\begin{document}

\addtobeamertemplate{block end}{}{\vspace*{2ex}} % White space under blocks
\addtobeamertemplate{block alerted end}{}{\vspace*{2ex}} % White space under highlighted (alert) blocks

\setlength{\belowcaptionskip}{2ex} % White space under figures
\setlength\belowdisplayshortskip{2ex} % White space under equations

\begin{frame}[t] % The whole poster is enclosed in one beamer frame

\begin{columns}[t] % The whole poster consists of three major columns, the second of which is split into two columns twice - the [t] option aligns each column's content to the top

\begin{column}{\sepwid}\end{column} % Empty spacer column

\begin{column}{\onecolwid} % The first column

%----------------------------------------------------------------------------------------
%	ABSTRACT
%----------------------------------------------------------------------------------------

\begin{block}{Abstract}
{
\small
Eigenvalues are ubiquitous in all branches of science and engineering. The
dominant eigenvalue algorithm reduces input matrices to almost diagonal and then
creates and chases bulges using parameters called shifts. Good parameter choices
can allow shrinking of the original problem allowing for a more aggressive
deflation on the subproblem. Such deflations greatly improve performance. We
apply automatic differentiation tools to achieve frequent deflations.
}
\end{block}

%----------------------------------------------------------------------------------------
%	EXISTING TECHNIQUES
%----------------------------------------------------------------------------------------

\begin{block}{Existing Techniques}
%
A general matrix is first reduced to Hessenberg~form, then the $\mathbf{Q}
\mathbf{R}$~\textbf{algorithm} is used to reduce the sub-diagonal entries to
zero iteratively:%
%
\begin{align*}
\vimage{ \begin{tikzonimage}[width=0.2\linewidth]{plots/original.eps}
\node[MyCaption] at (0.5, 1) {0 iter};
\end{tikzonimage} }
\!\!\vpointer\!\!
\vimage{ \begin{tikzonimage}[width=0.2\linewidth]{plots/qr1.eps}
\draw[-triangle 60, draw, line width=0.33ex, red] (0.1, 0.25) -- ++(0.2, 0.2);
\draw[-triangle 60, draw, line width=0.33ex, red] (0.25, 0.1) -- ++(0.2, 0.2);
\node[MyCaption] at (0.5, 1) {1 iter};
\end{tikzonimage} }
\!\!\vpointer\!\!
\vimage{ \begin{tikzonimage}[width=0.2\linewidth]{plots/qr2.eps}
\draw[-triangle 60, draw, line width=0.33ex, yellow] (0.1, 0.25) -- ++(0.2, 0.2);
\draw[-triangle 60, draw, line width=0.33ex, yellow] (0.25, 0.1) -- ++(0.2, 0.2);
\node[MyCaption] at (0.5, 1) {30};
\end{tikzonimage} }
\!\!\vpointer\!\!
\vimage{ \begin{tikzonimage}[width=0.2\linewidth]{plots/qrdone.eps}
\node[green] at (0.25, 0.25) {\Large $\checkmark$};
\node[MyCaption] at (0.5, 1) {$\cdots$};
\end{tikzonimage} }
\end{align*}
%\vspace{0.5ex}
%
Since the resulting matrix is strictly upper triangular, \textbf{the eigenvalues
are the diagonal entries}. The $Q R$~algorithm consists of repeated iterations
of:
%
\begin{align}
A_n     &\rightarrow Q_n R_n \label{eq:qr} \\
A_{n+1} &\leftarrow  R_n Q_n = Q_n^{-1} A_n Q_n \label{eq:rq}
\end{align}
%
As written, \eqref{eq:qr} and \eqref{eq:rq} require multiplication of large
matrices. For this reason, the $Q R$~algorithm is implemented as a loop of
similarity transformations.
%
\begin{align*}
\vimage{ \begin{tikzonimage}[width=0.2\linewidth]{plots/bulge1.eps}
\end{tikzonimage} }
\!\!\vpointer\!\!
\vimage{ \begin{tikzonimage}[width=0.2\linewidth]{plots/bulge2.eps}
\draw[-triangle 60, draw, line width=0.33ex, white] (0.05, 0.5) -- ++(0.25, -0.25);
\end{tikzonimage} }
\!\!\vpointer\!\!
\vimage{ \begin{tikzonimage}[width=0.2\linewidth]{plots/bulge3.eps}
\draw[-triangle 60, draw, line width=0.33ex, white] (0.2, 0.35) -- ++(0.25, -0.25);
\end{tikzonimage} }
\!\cdots\!
\vimage{ \begin{tikzonimage}[width=0.2\linewidth]{plots/bulgedone.eps}
\end{tikzonimage} }
\end{align*}
%
A \emph{bulge} is repeatedly \textbf{chased down the diagonal}. To speed
convergence \cite{watkins2007matrix}, the bulge is seeded with one or more
\emph{shifts}, which represent \textbf{parameter choices}.
% XXX SEPARATE
A second bulge can be
started from the bottom, and the two can meet in the middle.  A
\textbf{Schur decomposition is used to form \emph{spikes}} around this bulge.
Under the right conditions, the problem can be subdivided recursively into the
top-left and bottom-right quarters. This is shown below.
%
\begin{align*}
\vimage{ \begin{tikzonimage}[width=0.2\linewidth]{plots/bulges1.eps}
\end{tikzonimage} }
\!\!\vpointer\!\!
\vimage{ \begin{tikzonimage}[width=0.2\linewidth]{plots/bulges2.eps}
\end{tikzonimage} }
\!\!\vpointer\!\!
\vimage{ \begin{tikzonimage}[width=0.2\linewidth]{plots/bulgesspiked.eps}
\end{tikzonimage} }
\!\!\vpointer\!\!
\vimage{ \begin{tikzonimage}[width=0.2\linewidth]{plots/bulgesdeflated.eps}
\end{tikzonimage} }
\end{align*}

\vfill
\end{block}
\end{column} % End of the first column

\begin{column}{\sepwid}\end{column} % Empty spacer column

\begin{column}{\twocolwid} % Begin a column which is two columns wide (column 2)


%----------------------------------------------------------------------------------------
%	Proposed Techniques: Inducing Middle Deflations
%----------------------------------------------------------------------------------------
\begin{alertblock}{Proposed Technique: Inducing Recursive Middle Deflation}
%We expect middle deflation to be possible by chasing bulges from both corners
%simultaneously: \cite{BramanThesis} \cite{WatkinsBulgeExchanges}
%\cite{Giles2008} \cite{Struthers2014}
%
\begin{align*}
\vimage{ \begin{tikzonimage}[width=0.16\linewidth]{plots2/impstep002.eps} \end{tikzonimage} }
\!\!\vpointer\!\!
\vimage{ \begin{tikzonimage}[width=0.16\linewidth]{plots2/impstep006.eps} \end{tikzonimage} }
\!\!\vpointer\!\!
\vimage{ \begin{tikzonimage}[width=0.16\linewidth]{plots2/impstep008.eps} \end{tikzonimage} }
\!\!\vpointer\!\!
\vimage{ \begin{tikzonimage}[width=0.16\linewidth]{plots2/impstep009.eps} \end{tikzonimage} }
&
\!\!\vpointer\!\!
\vimage{ \begin{tikzonimage}[width=0.16\linewidth]{plots2/final.eps}
%\node[green] at (0.8, 0.7) {\huge $\checkmark$};
\draw[black,    line width=0.5ex]  (0,1)     rectangle (0.6,0.4);
\draw[green!90, line width=0.20ex] (0,1)     rectangle (0.6,0.4);
\draw[black,    line width=0.5ex]  (0.6,0.4) rectangle (1,0);
\draw[white!90, line width=0.20ex] (0.6,0.4) rectangle (1,0);
\end{tikzonimage} }
\\[1ex]
%
\vimage{ \begin{tikzonimage}[width=0.096\linewidth]{plots3/impstep002.eps}
\draw[black,    line width=0.5ex]  (0,1)     rectangle (1,0);
\draw[green!90, line width=0.20ex] (0,1)     rectangle (1,0);
\end{tikzonimage} }
\!\!\vpointer\!\!
\vimage{ \begin{tikzonimage}[width=0.096\linewidth]{plots3/impstep004.eps}
\draw[black,    line width=0.5ex]  (0,1)     rectangle (1,0);
\draw[green!90, line width=0.20ex] (0,1)     rectangle (1,0);
\end{tikzonimage} }
\!\!\vpointer\!\!
\vimage{ \begin{tikzonimage}[width=0.096\linewidth]{plots3/impstep005.eps}
\draw[black,    line width=0.5ex]  (0,1)     rectangle (1,0);
\draw[green!90, line width=0.20ex] (0,1)     rectangle (1,0);
\end{tikzonimage} }
&
\!\!\vpointer\!\!
\vimage{ \begin{tikzonimage}[width=0.096\linewidth]{plots3/final.eps}
%\node[green] at (0.8, 0.7) {\Large $\checkmark$};
\draw[black,    line width=0.5ex]  (0,1)     rectangle (1,0);
\draw[green!90, line width=0.20ex] (0,1)     rectangle (1,0);
\end{tikzonimage} }
%
\\[1ex]
%
\vimage{ \begin{tikzonimage}[width=0.064\linewidth]{plots4/impstep001.eps}
\draw[black,    line width=0.5ex]  (0,1)     rectangle (1,0);
\draw[white!90, line width=0.20ex] (0,1)     rectangle (1,0);
\end{tikzonimage} }
\!\!\vpointer\!\!
\vimage{ \begin{tikzonimage}[width=0.064\linewidth]{plots4/impstep002.eps}
\draw[black,    line width=0.5ex]  (0,1)     rectangle (1,0);
\draw[white!90, line width=0.20ex] (0,1)     rectangle (1,0);
\end{tikzonimage} }
\!\!\vpointer\!\!
\vimage{ \begin{tikzonimage}[width=0.064\linewidth]{plots4/impstep003.eps}
\draw[black,    line width=0.5ex]  (0,1)     rectangle (1,0);
\draw[white!90, line width=0.20ex] (0,1)     rectangle (1,0);
\end{tikzonimage} }
&
\!\!\vpointer\!\!
\hspace{0.096\textwidth}\hspace{0.5em} \vimage{ \begin{tikzonimage}[width=0.064\linewidth]{plots4/final.eps}
%\node[green] at (0.8, 0.75) {\large $\checkmark$};
\draw[black,    line width=0.5ex]  (0,1)     rectangle (1,0);
\draw[white!90, line width=0.20ex] (0,1)     rectangle (1,0);
\end{tikzonimage} }
\end{align*}
%
% position this absolutely on the page (ick)
\begin{tikzpicture}[remember picture, overlay]
\node [shift={(-30.5in,-13.25in)}, align=center, anchor=center] at (current page.north east) {
	\Large            \textbf{Subdivision is a} \\[1.5ex]
	\huge \color{red} \textbf{HUGE} \\[0.25ex]
	\Large            \textbf{performance win.}
};
\end{tikzpicture}

\end{alertblock}

\begin{columns}[t,totalwidth=\twocolwid] % Split up the two columns wide column

\begin{column}{\onecolwid}\vspace{-.6in} % The first column within column 2 (column 2.1)


%----------------------------------------------------------------------------------------
%	Schur Decomposition
%----------------------------------------------------------------------------------------
\begin{block}{Schur Decomposition}
%
The Schur decomposition allows us to create spikes. For an input matrix $A$, it
forms symmetric $Q$ and diagonal $T$ such that \textbf{$\mathbf{Q}$ diagonalizes
$\mathbf{A}$ into $\mathbf{T}$}, as 
%
\begin{align}
A = Q^\intercal T Q %\quad \therefore Q A Q^\intercal = T
\end{align}
%
Consider a submatrix $A_\text{sub}$ and compute $Q_\text{sub}$:
%
\begin{align*}
\vimage{ \begin{tikzonimage}[width=0.28\linewidth]{plots/schur_Az.eps}
\node[MyCaption] at (0.5, 1) {$A_\text{sub}$};
\end{tikzonimage} }
\vpointer
\vimage{ \begin{tikzonimage}[width=0.28\linewidth]{plots/schur_Qz.eps}
\node[MyCaption] at (0.5, 1) {$Q_\text{sub}$};
\end{tikzonimage} }
\vpointer
\vimage{ \begin{tikzonimage}[width=0.28\linewidth]{plots/schur_Q.eps}
\node[MyCaption] at (0.5, 1) {$Q$};
\end{tikzonimage} }
\end{align*}
%
We can then create $Q$ by embedding  $Q_\text{sub}$ into the identity.
Applying this to $A$ creates zeros on the subdiagonal in the neighborhood of the
decomposition, and forms spikes.
%
\begin{align*}
\vimage{ \begin{tikzonimage}[width=0.28\linewidth]{plots/schur_Ahess.eps}
\node[MyCaption] at (0.5, 1) {$A$};
\end{tikzonimage} }
\vpointer
\vimage{ \begin{tikzonimage}[width=0.28\linewidth]{plots/schur_final.eps}
\node[MyCaption] at (0.5, 1) {$Q^\intercal A Q$};
\end{tikzonimage} }
\end{align*}
%
This matrix is then reduced to Hessenberg form again.
If half of the values on the spike are near zero, then
the Hessenberg form has a zero on the subdiagonal which
allows for matrix deflation (see above box).
%
\end{block}


%----------------------------------------------------------------------------------------

\end{column} % End of column 2.1

\begin{column}{\onecolwid}\vspace{-.6in} % The second column within column 2 (column 2.2)


%----------------------------------------------------------------------------------------
%	Automatic Differentiation
%----------------------------------------------------------------------------------------
\begin{block}{Achieving Deflation}

In practice, middle deflations are improbable because they rely on the spikes
having many zeros. The values on the spike depend on the shifts. Since we can
control these shifts, the problem of creating a middle deflation can be phrased
as
%
\begin{equation}
  \min_{\textrm{shifts}} \left( \Sigma 
    \left| \textrm{selected spike values}%I'm not sure if 'selected' is necesary. True, but maybe cluttersome
  \right| \right)
\end{equation}
%
Minimizing this implies that the spikes have many zeros, allowing for middle
deflation. Selecting a single element and tracing its behavior would cause it
to sketch out a curve as shown below:

\begin{align*}
\vimage{ \begin{tikzonimage}[width=0.40\linewidth]{plots2/impstep009.eps}
		\draw[red, ultra thick] (0.325, 0.375) circle (1ex) node[left=0.5em,
		rounded corners, white, fill=black, opacity=0.95, text opacity=1.0] {$x$};
\end{tikzonimage} }
\!\!\vpointer\!\!
\vimage{ \begin{tikzpicture}
\begin{axis}[
		xlabel={(shift choices)}, ylabel=$x$,
	width=0.45\linewidth,
	axis x line=center,
	axis y line=center,
	xtick=\empty, ytick=\empty,
	no markers,
	x label style={at={(axis description cs:0.5,0.0)},anchor=north},
]
\addplot[red] {(x-2)^2 - (x-2) - 0};
\end{axis}
\end{tikzpicture} }
\end{align*}

\begin{align*}
\end{align*}
%
Many effective methods for minimizing functions rely upon derivatives (e.g.
Newton's Method). \textbf{Automatic differentiation} (AD) methods allow for
differentiation of matrix operations \cite{Giles2008} and the schur
decomposition \cite{Struthers2014}. Thus by using AD tecniques we can utilize
existing minimization algorithms to create deflatable matrices.

\end{block}


%----------------------------------------------------------------------------------------
%	FUTURE WORK
%----------------------------------------------------------------------------------------
%\begin{block}{Future Work}
%In this section, cover:
%\begin{itemize}
%\item What are we going to do next? Very much depends on results that have not
%yet been teased out.
%\end{itemize}
%\end{block}

%----------------------------------------------------------------------------------------

\end{column} % End of column 2.2

\end{columns} % End of the split of column 2 - any content after this will now take up 2 columns width

%----------------------------------------------------------------------------------------

%\begin{columns}[t,totalwidth=\twocolwid] % Split up the two columns wide column again
%
%\begin{column}{\onecolwid} % The first column within column 2 (column 2.1)
%
%\end{column} % End of column 2.1
%
%\begin{column}{\onecolwid} % The second column within column 2 (column 2.2)
%
%%----------------------------------------------------------------------------------------
%
%\end{column} % End of column 2.2
%
%\end{columns} % End of the split of column 2

%----------------------------------------------------------------------------------------
%	PERFORMANCE IMPROVEMENT
%----------------------------------------------------------------------------------------

%\begin{alertblock}{Performance Improvement}
%
%Our algorithm outperforms XXX by XXXX. We show this by using a MATLAB package
%which counts FLOPs which shows that our algorithm performs fewer operations when
%compared to traditional QR on the same matrix.
%
%\begin{figure}[h]
%(pretty pictures here)
%\end{figure}
%
%Other text TBD
%
%Note that this whole poster is now in English (no Latin :) ).
%
%There is a very cool-looking package that might work with MATLAB for counting
%FLOPS:
%\url{http://research.microsoft.com/en-us/um/people/minka/software/lightspeed/}
%
%\end{alertblock}

\end{column} % End of the second column

\begin{column}{\sepwid}\end{column} % Empty spacer column

\begin{column}{\onecolwid} % The third column

%----------------------------------------------------------------------------------------
%	CONCLUSIONS
%----------------------------------------------------------------------------------------

\begin{block}{Conclusions}

We have presented an algorithm which aggresively targets
middle deflation as opposed to existing algorithms.
Furthermore, its recursive nature exposes parralizability
which could lead 

With the proper implementation, this will provide a
large performance improvement over traditional algorithms in use today.
Additionally, this is an enabling technique for extremely large eigenvalue
%problems. Our preliminary results indicate that a performance gain of XXX can be
%expected from proper implementations of this technique. This was determined by
%naively implementing each algorithm in MATLAB and comparing FLOP counts using
%each algorithm for the same inputs.

\end{block}

%----------------------------------------------------------------------------------------
%	FUTURE WORK
%----------------------------------------------------------------------------------------

\begin{block}{Future Work}

\begin{itemize}
  \item Incorporate AD methods into minimization algorithms
  \item Characterize the performance of the algorithm
  \item Optimize performance
  \begin{itemize}
    \item Time
    \item Memory
  \end{itemize}
\end{itemize}

\href{http://octave.org/}{GNU Octave}.

\end{block}

%----------------------------------------------------------------------------------------
%	REFERENCES
%----------------------------------------------------------------------------------------

\begin{block}{References}

%\nocite{*} % Insert publications even if they are not cited in the poster
\small{\bibliographystyle{siam}
\bibliography{papers}\vspace{0.75in}}

\end{block}

%----------------------------------------------------------------------------------------
%	ACKNOWLEDGEMENTS
%----------------------------------------------------------------------------------------

%\setbeamercolor{block title}{fg=red,bg=white} % Change the block title color

\begin{block}{Acknowledgements}

The authors thank \textbf{Dr.~Allan~Struthers} (in the Mathematics department)
for his guidance and creation of AD tools for the Schur decomposition.

\end{block}

%----------------------------------------------------------------------------------------
%	CONTACT INFORMATION
%----------------------------------------------------------------------------------------

\setbeamercolor{block alerted title}{fg=black,bg=norange} % Change the alert block title colors
\setbeamercolor{block alerted body}{fg=black,bg=white} % Change the alert block body colors

\begin{block}{Contact Information}

\begin{itemize}
%\item Web: \href{http://www.university.edu/smithlab}{http://www.university.edu/smithlab}
\item Email: \href{mailto:ahirzel@mtu.edu}{\texttt{ahirzel@mtu.edu}}
\item Phone: +1 (906) 231 0866
\end{itemize}

\end{block}

\begin{center}
% minimum size is 1 3/4"
\includegraphics[width=0.4\linewidth]{logo_withfullname_LG.eps}
\end{center}

%----------------------------------------------------------------------------------------

\end{column} % End of the third column

\end{columns} % End of all the columns in the poster

\end{frame} % End of the enclosing frame

\end{document}
